\subsection{Experimental Setup}

This section describes the experiments which are conducted to evaluate the observability model's performance, and determine the impact of the VBEE library on SLAM systems. These experiments are defined here, and the results are presented in Section \ref{sec:experimental_results}.

\subsubsection{Parameter Optimization}

Before experimentation can take place on the SLAM system, the parameters of the VBEE library must be tuned to optimize performance. Parameters are split between the VBEE framework, the observability model, and the existence probability estimator.  These components are tightly coupled, as each provides feedback to the others, making optimization difficult due to the high number of parameters. For this reason, parameter optimization takes place in three phases. The first phase freezes the parameters of the VBEE framework and the existence probability estimator to reasonable values, and optimizes the observability model's parameters. The second phase freezes the observability model's parameters to the obtained values, optimizing the VBEE and EPE parameters. The final phase tunes all parameters simultaneously, ideally achieving optimal performance for the system.

\paragraph{Phase 1 - Observability Model Tuning}

The observability model has parameters $\params{\knn} = \{k, n, \athreshold, \feedbackthreshold\}$, which are tuned to minimize the following objectives:
\begin{itemize}
    \item $n_{\text{converge}}$: The number of observations necessary for the VBEE system to converge to a stable estimate of $\pe$
    \item $\epsilon_{\text{converge}}$: The observability model's error after $n_{\text{converge}}$ observations
    \item $n_{\text{adapt}}$: The number of observations necessary for the system to re-converge to a stable estimate of $\pe$ after changing the point's observability
    \item $\epsilon_{\text{adapt}}$: The observability model's error after $n_{\text{adapt}}$ observations
    \item $t_{update}$: The average time in ms needed to perform one estimation and one update
    \item size: The total storage in bytes needed to serialize one map point's observability
\end{itemize}

In this first phase, the $\vbee$ and $\epe$ parameters are locked to the following values:
\begin{align*}
    \params{\epe}^* &= \{\falseneg = 0.1, \falsepos = 0.001\}\\
    \params{\vbee}^* &= \{\estpe_0 = 0.9, \damping = 0.1, \observability_0 = 0.5, \observabilityupdate = 0.5\}
\end{align*}
The observability model parameters $\params{\knn}$ are optimized to minimize the objectives using the objective functions described below.

The function $\operatorname{ObjectiveNConverge}$ shown in Algorithm \ref{alg:obj_n_converge} determines the average number of observations needed for the system to converge to a stable $\estpe$ for each observability scenario in the dataset. The system is considered converged when the standard deviation over the last $w = 100$ samples drops below $\sigma_{\text{converge}} = 0.2$. If the system does not converge within $m_{\text{max}} = 1000$ observations, the test is considered failed, and the optimization is penalized accordingly.

\begin{singlespace}
    \begin{algorithm}[H]
        \caption[Objective Function: $n_{converge}$]{ObjectiveNConverge}
        \label{alg:obj_n_converge}
        \begin{algorithmic}[1]
            \Require{$vbee$: VBEE instance initialized with $\params{\vbee}^*,\params{\knn}, \params{\epe}^*$}
            \Require{$observability\_scenarios$: Observability Scenario Dataset}
            \Function{ObjectiveNConverge}{}
            \State{$all\_n\_convergences \gets [\;]$}
            \For{$os$ in $observability\_scenarios$}
                \State{$\texttt{reset } vbee$} \Comment{Reset VBEE to initial state for each observability scenario}
                \State{$window \gets \operatorname{Queue()}$}
                \State{$n\_converge \gets -1$}
                \For{$i = 0$ to $ i = m_{max}$}
                    \State{$v \gets \texttt{random viewpoint}$}
                    \State{$new\_pE \gets vbee.\update{v, scenario.\operatorname{Query}(v)}$}
                    \State{$\texttt{push } new\_pE$ onto $window$}
                    \If{$window \texttt{ size}$ $> w$}
                        \State{$\texttt{pop }window$}
                        \If{$\texttt{std dev}$ of $window < \sigma_{converge}$}
                            \State{$n\_converge \gets i$}
                            \State{break}
                        \EndIf
                    \EndIf
                \EndFor
                \If{$n\_converge = -1$}
                    \State{$\texttt{append } 2 * m_{max}$ to $all\_n\_convergences$}
                \Else
                    \State{$\texttt{append } i$ to $all\_n\_convergences$}
                \EndIf
            \EndFor
            \State{\Return{$\texttt{average } all\_n\_convergences$}}
            \EndFunction
        \end{algorithmic}
    \end{algorithm}
\end{singlespace}

The function $\operatorname{ObjectiveEConverge}$ shown in Algorithm \ref{alg:obj_e_converge} determines average error between the observability model and the observability scenario after the system has converged across all observability scenarios. Model error is calculated as the sum of differences between the model's estimate of $\pe$, and the true observability from the observability scenarios dataset over $j = 1000$ random viewpoints. Once again, if the system does not converge to a stable probability estimate, the system is penalized.

\begin{singlespace}
    \begin{algorithm}[H]
        \caption[Objective Function: $\epsilon_{converge}$]{ObjectiveEConverge}
        \label{alg:obj_e_converge}
        \begin{algorithmic}
            \Require{$vbee$: VBEE instance initialized with $\params{\vbee}^*,\params{\knn}, \params{\epe}^*$}
            \Require{$observability\_scenarios$: Observability Scenario Dataset}
            \Function{ObjectiveEConverge}{}
            \State{$all\_e\_convergences \gets [\;]$}
            \For{$os$ in $observability\_scenarios$}
                \State{$\texttt{reset } vbee$} \Comment{Reset VBEE to initial state for each observability scenario}
                \State{$error \gets 0$}
                \State{$window \gets \operatorname{Queue()}$}
                \For{$m_{max}$ times}
                    \State{}
                    \State{[\textit{Same as Algorithm \ref{alg:obj_n_converge}, lines} 8-13]}
                    \State{}
                    \If{$window \texttt{ size}$ $> w$}
                        \State{$\texttt{pop }window$}
                        \If{$\texttt{std dev}$ of $window < \sigma_{converge}$}
                            \For{$j$ times}
                                \State{$v \gets \texttt{random viewpoint}$}
                                \State{$true\_obs \gets os.\operatorname{Query}(v)$}
                                \State{$model\_obs \gets vbee.model.\estimate{v}$}
                                \State{$error \gets error + |true\_obs - model\_obs|$}
                            \EndFor
                            \State{break}
                        \EndIf
                    \EndIf
                \EndFor
                \If{$error = 0$}
                    \State{$\texttt{append } 1000$ to $all\_e\_convergences$}
                \Else
                    \State{$\texttt{append } error$ to $all\_e\_convergences$}
                \EndIf
            \EndFor
            \State{\Return{$\texttt{average } all\_e\_convergences$}}
            \EndFunction
        \end{algorithmic}
    \end{algorithm}
\end{singlespace}

The functions ObjectiveNAdapt and ObjectiveEAdapt are implemented similarly to ObjectiveNConverge and ObjectiveEConverge respectively, but perform two convergence loops for each observability scenario. Between the two convergence loops, the active observability scenario is replaced with a predetermined partner scenario, selected to be significantly different from the original, but with some observability overlap.

The function ObjectiveTUpdate, shown in Algorithm \ref{alg:obj_t_update} determines the average time to perform VBEE update for 1000 random observations, averaged for every observability scenario.

\begin{singlespace}
    \begin{algorithm}[H]
        \caption[Objective Function: $t_{update}$]{ObjectiveTUpdate}
        \label{alg:obj_t_update}
        \begin{algorithmic}
            \Require{$vbee$: VBEE instance initialized with $\params{\vbee}^*,\params{\knn}, \params{\epe}^*$}
            \Require{$observability\_scenarios$: Observability Scenario Dataset}
            \Function{ObjectiveTUpdate}{}
            \State{$all\_t\_updates \gets [\;]$}
            \For{$os$ in $observability\_scenarios$}
                \State{$\texttt{reset } vbee$} \Comment{Reset VBEE to initial state for each observability scenario}
                \State{$avg\_t\_update \gets 0$}
                \For{1000 times}
                    \State{$v \gets \texttt{random viewpoint}$}
                    \State{$t_0 \gets \texttt{current time in ms}$}
                    \State{$new\_pE \gets vbee.\update{v, scenario.\operatorname{Query}(v)}$}
                    \State{$t_1 \gets \texttt{current time in ms}$}
                    \State{$avg\_t\_update \gets avg\_t\_update + (t_1 - t_0$)}
                \EndFor
                \State{$\texttt{append } (avg\_t\_update / 1000)$ to $all\_t\_updates$}
            \EndFor
            \State{\Return{$\texttt{average } all\_t\_updates$}}
            \EndFunction
        \end{algorithmic}
    \end{algorithm}
\end{singlespace}

The function ObjectiveSize calculates the total size in bytes necessary to store all information necessary to reconstruct a serialized map point VBEE system in the same state as it was when it was serialized. This reconstruction requires:
\begin{itemize}
    \item All parameter vectors $\params{\vbee},\params{\knn}, \text{ and } \params{\epe}$
    \item The current estimated existence probability $\estpe$
    \item The current observability $\observability$
    \item The list of observations from the observability model
\end{itemize}
The only item in this list with a size dependent on the parameter vectors is the observations list, which has a maximum capacity of $n$, provided in $\params{\knn}$. All other sizes are not dependent on the input parameters, and can therefore be left out of the objective function. The serialization of an observation requires three floating point values for the x, y, and z coordinates in the map point's frame, and a boolean value for whether the observation is positive or negative. If a float is assumed to be three bytes in size, and the boolean value will also require one byte, then ObjectiveSize can be implemented to simply return $4 * n$.

The final multi-objective function is shown in Algorithm \ref{alg:obj_obs_model}. This function returns an objective vector of six values, each related to one of the performance criteria defined previously. At this stage, the system is able to undergo multi-objective optimization to find values which minimize each of the six objectives.

\begin{singlespace}
    \begin{algorithm}[H]
        \caption{Observability Model Objective Function}
        \label{alg:obj_obs_model}
        \begin{algorithmic}
            \Require{$\params{\vbee}^*, \params{\epe}^*$: Static VBEE and EPE parameters}
            \Function{ObjectiveObservabilityModel}{$\params{\knn}$}
            \State{$vbee \gets \vbee.\initialize{\params{\vbee}^*, \params{\knn}, \params{\epe}^*}$}
            \State{$n_{converge} \gets \operatorname{ObjectiveNConverge}()$}
            \State{$\texttt{reset } vbee$} \Comment{Reset VBEE to initial state between objectives}
            \State{$\epsilon_{converge} \gets \operatorname{ObjectiveEConverge}()$}
            \State{$\texttt{reset } vbee$}
            \State{$n_{adapt} \gets \operatorname{ObjectiveNAdapt}()$}
            \State{$\texttt{reset } vbee$}
            \State{$\epsilon_{adapt} \gets \operatorname{ObjectiveEAdapt}()$}
            \State{$\texttt{reset } vbee$}
            \State{$t_{update} \gets \operatorname{ObjectiveTUpdate}()$}
            \State{$\texttt{reset } vbee$}
            \State{$\text{size} \gets \operatorname{ObjectiveSize}()$}
            \State{\Return{$(n_{converge},\epsilon_{converge},n_{adapt},\epsilon_{adapt},t_{update},\text{size})$}}
            \EndFunction
        \end{algorithmic}
    \end{algorithm}
\end{singlespace}

The python library Optuna \cite{akibaOptunaNextgenerationHyperparameter2019} was selected to perform the multi-objective optimization of the system's parameters. Optuna is a hyperparameter optimization framework which allows for the discovery of Pareto optimal \cite{censorParetoOptimalityMultiobjective1977} parameter sets. 